\section{What is said and what settles}
\label{sec:tactics}

The prompt invites agents to ``negotiate, bluff, deceive, cooperate, intimidate'' in public, so a public bluff is permitted conduct here, not a violation of instructions. The incidents above already show selective disclosure (G13) and a withdrawn assurance (G18). Two more show how claims and terms come apart. In one a bluff settles; in the other the recipient rejects the claim and accepts the deal.

\TranscriptTitle{g7}{A ``final'' price below a private limit}
Sol asks \$625 for North Carolina, which Grok wants for its greens. Grok counters with a public limit.
\TranscriptQuote{g7_d42_public}
Its private note on the same response allows more.
\TranscriptQuote{g7_d42_private}
\TranscriptQuote{g7_d43_public}
\TranscriptFinding{The deed sells for \$600. Sol's note prefers reinvesting the cash to holding out for \$25. Grok completes green at turn 95 by buying Pennsylvania Avenue at auction for \$171 but never builds on it. At turn 110 it is eliminated owing Sol \$1,400, and North Carolina returns to Sol inside a complete green group. Sol builds on every green two turns later.}
\TranscriptReading{The public limit is below the private one, and the bluff settles at the stated price. It contrasts with G2, where Opus's ``final offer'' was refused. The \$25 at stake is small, and nothing in the record isolates the bluff's effect on the price. The later estate returns the same deed to its seller, now as part of a complete group.}
\TranscriptSource{G7 d41--43, standard first attempts; auction event 1680, estate events 1946--1954, build events 1971--1973. The \$600 message appears in Sol's next request.}

\TranscriptTitle{g19}{A false claim corrected, and the deal accepted}
Grok's received call offers Virginia Avenue and \$150 for States Avenue, with this public message.
\TranscriptQuote{g19_d135_false_public}
Both deeds are pink. Grok's note on the same response has already noticed the error.
\TranscriptQuote{g19-abort}
Its assistant content names \End, but the provider returned the proposal, and the engine applied it. Sol replies.
\TranscriptQuote{g19_d136_correction_accept}
\TranscriptFinding{The swap and the \$150 execute and complete no group. Sol's note calls it a ``strictly favorable swap.''}
\TranscriptReading{Sol separates the terms from the explanation. It corrects the false colour claim in public and takes the cash and deed on their merits. Acceptance of an offer is therefore not acceptance of the proposer's account of it. The proposal is one of Grok's conflicting calls, so it does not show what Grok intended.}
\TranscriptSource{G19 d135--136, standard first attempts. The false claim appears in Sol's delivered request.}

The agents' messages carry claims about cash, probabilities, colours and resale, and several of those claims are wrong or incomplete. In the incidents, recipients' replies turn on the terms and holdings, which the prompt states exactly. Sol corrects Grok's colour claim in G19, and Opus does not endorse Grok's cash limit in G13. We did not measure how often a claim changes a response. The population contrasts in Section~\ref{sec:bargaining} follow structural facts about each exchange that every recipient could see, which is consistent with this reading of the incidents.
